\section{Arquitectura del sistema}
\label{sec:arquitectura}

Esta sección documenta la arquitectura resultante tras la Entrega 4 en los tres niveles del
modelo C4~\cite{brown2018c4model}, y resume las decisiones registradas como \emph{Architectural
Decision Records} (ADR) que la sustentan.

\subsection{Vista de contexto (C4 Nivel 1)}

La Figura~\ref{fig:c4-nivel1} sitúa el sistema como una sola caja frente a sus tres tipos de
usuario. No hay sistemas externos reales que integrar en este nivel: lo que podría parecer una
integración externa —notificaciones multicanal, telemetría de equipos del abonado— está simulado
dentro del propio sistema y documentado como tal en cada ADR correspondiente (ver
\texttt{docs/diagrams/c4-nivel1-contexto.md} para el detalle de por qué).

\begin{figure}[H]
\centering
\includegraphics[width=0.85\textwidth]{../diagrams/c4-nivel1-contexto.png}
\caption{C4 Nivel 1 — contexto del sistema.}
\label{fig:c4-nivel1}
\end{figure}

\subsection{Vista de contenedores (C4 Nivel 2)}

El diagrama de contenedores de la Entrega 3 (\texttt{docs/diagrams/c4-nivel2-contenedores.md})
se actualiza en esta entrega con las piezas nuevas: las dos
aplicaciones cliente (\texttt{apps/web}, \texttt{apps/mobile}) ya no hablan directamente con cada
microservicio por su puerto, sino que pasan por un único \texttt{api-gateway} (Spring Cloud
Gateway) que enruta por prefijo de \emph{path} sin reescribir el cuerpo de la petición. Se agrega
\texttt{telemetry-service} (PE-U1, Sección~\ref{sec:pe-u1-correl}) como séptimo microservicio, y
el stack de observabilidad (OpenTelemetry Collector, Tempo, Prometheus, Grafana, cAdvisor) como
un grupo de contenedores transversal. La Figura~\ref{fig:c4-nivel2} muestra el diagrama completo.

\begin{figure}[H]
\centering
\includegraphics[width=\textwidth]{../diagrams/c4-nivel2-contenedores.png}
\caption{C4 Nivel 2 — contenedores del sistema, Entrega 4.}
\label{fig:c4-nivel2}
\end{figure}

\begin{table}[H]
\centering
\small
\caption{Contenedores del sistema tras la Entrega 4.}
\label{tab:contenedores-e4}
\begin{tabular}{@{}p{3.1cm}p{2.6cm}p{7.8cm}@{}}
\toprule
\textbf{Contenedor} & \textbf{Tecnología} & \textbf{Responsabilidad} \\
\midrule
\texttt{apps/web} & React 18 + TypeScript & SPA de consola para clientes/técnicos/administradores \\
\texttt{apps/mobile} & Kotlin + Jetpack Compose & Cliente de campo para técnicos (APK) \\
\texttt{api-gateway} & Spring Cloud Gateway & Único punto de entrada, enrutamiento por prefijo \\
\texttt{auth-service} & Java 21 + Spring Boot & JWT, RBAC, refresh tokens \\
\texttt{ticket-service} & Java 21 + Spring Boot & CRUD de tickets, arquitectura hexagonal (Sección~\ref{sec:hexagonal-e4}) \\
\texttt{ai-service} & Python + FastAPI & Clasificación asíncrona por reglas \\
\texttt{notification-service} & Node.js + Express & Notificaciones multicanal simuladas \\
\texttt{report-service} & Java 21 + Spring Boot & Modelo de lectura CQRS, reportes \\
\texttt{telemetry-service} & Java 21, sockets TCP + gRPC & Canal PE-U1: equipos del abonado + latidos de nodo, reloj de Lamport (Sección~\ref{sec:pe-u1-correl}) \\
CockroachDB & Cluster 3 nodos & Persistencia distribuida (auth\_db/ticket\_db/report\_db) \\
Kafka + MongoDB & --- & Saga por coreografía, clasificaciones y notificaciones \\
OTel Collector, Tempo, Prometheus, Grafana, cAdvisor & --- & Observabilidad (Sección~\ref{sec:observabilidad}) \\
\bottomrule
\end{tabular}
\end{table}

\subsection{Arquitectura hexagonal de \texttt{ticket-service}}
\label{sec:hexagonal-e4}

El Módulo A de esta entrega exigió refactorizar el microservicio principal en cuatro capas con
una regla de dependencia estricta: las capas externas dependen de las internas, nunca al revés.
Antes del refactor, una única clase \texttt{TicketService} mezclaba acceso a datos vía
\texttt{JpaRepository} inyectado directamente, la lógica de creación de tickets en línea, dos
copias distintas del cálculo de SLA, y los tres casos de uso mutables como métodos sin forma
uniforme de interceptarlos.

\begin{table}[H]
\centering
\small
\caption{Las cuatro capas de \texttt{ticket-service} tras el refactor.}
\label{tab:capas-e4}
\begin{tabular}{@{}p{2.6cm}p{5.8cm}p{5.1cm}@{}}
\toprule
\textbf{Capa} & \textbf{Contenido} & \textbf{Regla} \\
\midrule
\texttt{presentation} & \texttt{TicketController}, DTOs & Traduce HTTP $\leftrightarrow$ comandos; sin lógica de negocio \\
\texttt{application} & \texttt{*CommandHandler} & Orquesta el caso de uso; no conoce HTTP ni JPA \\
\texttt{domain} & \texttt{Ticket}, puertos, políticas & Cero anotaciones Spring/JPA, verificado
  con una prueba ArchUnit que rompe la construcción si deja de serlo; se prueba con JUnit puro \\
\texttt{infrastructure} & Adaptadores, filtros, métricas, \emph{scheduler} & Implementa los puertos del dominio contra CockroachDB, Kafka y HTTP \\
\bottomrule
\end{tabular}
\end{table}

\subsection{Seis patrones GoF}

Se implementaron seis patrones del catálogo de Gamma et al.~\cite{gamma1994patterns} —la unión
del mínimo exigido por el Módulo A (Repository, Factory Method, Strategy) y del conjunto
asignado específicamente al equipo ACC en la Tabla~1 de la guía de la entrega (Command, Chain of
Responsibility, Observer, Repository, Strategy)—, cada uno documentado en detalle como
\href{run:../adr/0005-patrones-gof.md}{ADR-0005} con el problema de diseño real que resolvía
antes del refactor:

\begin{enumerate}[nosep]
  \item \textbf{Repository} --- desacopla el dominio de \texttt{JpaRepository} (Spring Data) vía
        un puerto (\texttt{TicketRepository}, en el paquete \texttt{domain}) implementado por un
        adaptador de infraestructura.
  \item \textbf{Factory Method} --- centraliza la creación de un \texttt{Ticket} nuevo
        (\texttt{TicketFactory.crearNuevo}), antes duplicada en línea.
  \item \textbf{Strategy} --- dos políticas de SLA intercambiables (\texttt{DefaultSlaPolicy},
        \texttt{ClassifiedSlaPolicy}) que antes vivían hardcodeadas en dos lugares distintos.
  \item \textbf{Command} --- cada acción mutable (crear, cambiar estado, asignar técnico) es un
        objeto inmutable ejecutado por un \emph{handler} dedicado, desacoplando al controlador de
        cómo se ejecuta la acción.
  \item \textbf{Chain of Responsibility} --- dos eslabones de evaluación de escalado
        (\texttt{SlaBreachedEscalationHandler}, \texttt{StaleCriticalEscalationHandler}) que
        implementan una funcionalidad que antes no existía en absoluto: el estado
        \texttt{ESCALADO} era inalcanzable desde la Entrega 3.
  \item \textbf{Observer} --- tres observadores del evento de escalado (\emph{logging}, métrica de
        negocio, publicación en Kafka), agregables sin modificar el sujeto que los notifica.
\end{enumerate}

Cada patrón reemplaza una duplicación o acoplamiento verificable contra el historial de
\texttt{TicketService.java} previo a esta entrega, no una envoltura artificial agregada solo para
completar el conteo mínimo.

\subsection{Canal de telemetría PE-U1 y correlación de tickets en Incidencias}
\label{sec:pe-u1-correl}

Se construyó \texttt{telemetry-service}: un servidor de sockets TCP que recibe reportes
periódicos de equipos del abonado y latidos de los 3 nodos del clúster, les asigna un
\emph{timestamp} de orden causal mediante un reloj lógico de Lamport~\cite{lamport1978time}, y
expone el registro resultante por gRPC (\texttt{GetEventosPorZona}), ordenado por ese
\emph{timestamp} y no por orden de llegada (ADR-0007). Dos clientes simuladores, cada uno con su
propio reloj de Lamport independiente del servidor, demuestran la regla clásica
$\text{max}(local, recibido) + 1$ entre procesos reales.

Sobre ese canal se construyó el mecanismo de correlación de tickets en \emph{Incidencias}: un
ticket nuevo puede agruparse con otros de la misma zona si la estrategia activa (variable de
entorno \texttt{CORREL}, patrón Strategy) lo decide. Se implementaron tres estrategias
intercambiables (ADR-0008), seleccionadas por Spring inyectando un
\texttt{Map<String, CorrelationStrategy>} y resolviendo la activa por el nombre exacto del
\emph{bean}:

\begin{itemize}[nosep]
  \item \textbf{\texttt{c0}} (línea base): cada ticket abre su propia incidencia.
  \item \textbf{\texttt{c1}}: se une a la incidencia abierta más reciente de la misma zona dentro
        de una ventana deslizante, ciega a la telemetría.
  \item \textbf{\texttt{c2}}: mismo criterio de \texttt{c1}, más una consulta real por gRPC a
        \texttt{telemetry-service}; solo agrupa si hay evidencia real de telemetría en esa
        ventana. Si el canal no responde, el ticket queda aislado en su propia incidencia --
        nunca bloquea ni revierte la creación del ticket (mismo principio de fallo abierto que
        ya aplica el \emph{publish} de Kafka, ADR-0004).
\end{itemize}

Ninguna de las tres estrategias está diseñada para resolver el caso de dos averías reales
golpeando la misma zona en la misma ventana -- ambas pueden fundirlas en una sola incidencia. Es
intencional: el objetivo de un escenario con dos averías simultáneas en el protocolo experimental
es \emph{revelar} ese error, no demostrar que no ocurre. Se construyó además un inyector de
averías (\texttt{experimentos/inyector\_averias.py}) que provoca una avería simulada real
--crea tickets reales vía la API y envía telemetría real al canal-- y deja constancia exacta de
a quién afectó en \texttt{verdad\_campo.csv}, precisamente porque en un ISP real nadie sabe con
certeza qué abonados estaban afectados por una avería y aquí sí se conoce con exactitud. Una
primera corrida manual del inyector junto con el script de análisis
(\texttt{experimentos/analizar\_correlacion.py}) confirmó el mecanismo funcionando de punta a
punta con datos reales: 100\,\% de precisión y exhaustividad en el agrupamiento bajo \texttt{c2}
con telemetría real corroborando la avería. La batería estadística completa del protocolo (10
repeticiones $\times$ 4 escenarios $\times$ 3 modos) es trabajo de cierre pendiente, no incluido
en esta corrida de verificación.

\subsection{Postura de consistencia distribuida}

La Entrega 3 estableció (ADR-0004) una postura CAP/PACELC~\cite{brewer2012cap,abadi2012consistency}
mixta por agregado de dominio: la creación de tickets prioriza disponibilidad (AP), mientras que
los datos de SLA y facturación exigen consistencia fuerte (CP). CockroachDB no ofrece un modo
``eventual configurable por tabla —su consistencia serializable por defecto, lograda mediante
consenso Raft~\cite{ongaro2014raft}, cubre el caso CP sin configuración adicional; la postura AP
se aproxima evitando reintentos automáticos silenciosos en las operaciones críticas. Esta postura
sigue vigente sin cambios en la Entrega 4: el refactor a arquitectura hexagonal reorganizó
\emph{dónde} vive la lógica de persistencia, no \emph{cómo} se comporta frente a particiones de
red.
